% -- Active: 1770635456760@@127.0.0.1@3306
%%%%%%%%%%%%%%%%%%%%%%%%%%%%%%%%%%%%%%%%%%%%%%%%%%%%%%%%%%%%%%%
%
% SCHEMAT HYUNDAI COUPE MULTIGAUGE (ESP32)
% Poprawiona wersja v5.0 - Pełna strona A4 Landscape
%
%%%%%%%%%%%%%%%%%%%%%%%%%%%%%%%%%%%%%%%%%%%%%%%%%%%%%%%%%%%%%%%
\documentclass{article}
\usepackage{geometry}
% Ustawiamy szeroką stronę (Landscape)
\geometry{landscape, a4paper, margin=1cm} 
\usepackage[european, straightvoltages]{circuitikz}
\usepackage{tikz}
\usepackage[T1]{fontenc}
\usepackage[utf8]{inputenc}

% Konfiguracja stylów (Grubsze linie, czytelne czcionki)
\ctikzset{
    bipoles/thickness=1.2,   % Grubsze elementy
    font=\sffamily\small     % Czcionka bezszeryfowa
}

\begin{document}

\begin{center}
\textbf{\LARGE Schemat Elektryczny: Hyundai Coupe Multigauge (ESP32 + CAN)}\\[0.5cm]
\end{center}

\centering
% Skalowanie schematu, aby wypełnił stronę
\begin{circuitikz}[scale=0.9, transform shape]

% ==================================================================
% 1. WTYCZKA SAMOCHODU (M15-B) - LEWA STRONA
% ==================================================================
\draw (-9, 8) node[draw, thick, fill=lightgray!20, minimum width=3.5cm, minimum height=16cm, align=center, rounded corners] (CONN) {\textbf{WTYCZKA}\\\textbf{M15-B}\\\textit{(Samochód)}};

% Definicja pinów (Współrzędne przy krawędzi wtyczki)
\coordinate (P_IGN) at (-7.25, 7);   \node[left] at (-7.3, 7) {Pin 3: IGN (+12V)};
\coordinate (P_GND) at (-7.25, 6);   \node[left] at (-7.3, 6) {Pin 12: GND};
\coordinate (P_BAT) at (-7.25, 3);   \node[left] at (-7.3, 3) {Pin 6: BAT (+12V)};
\coordinate (P_ILL) at (-7.25, 0);   \node[left] at (-7.3, 0) {Pin 9: ILL (Światła)};
\coordinate (P_SPD) at (-7.25, -2);  \node[left] at (-7.3, -2) {Pin 4: SPEED};
\coordinate (P_INJ) at (-7.25, -4);  \node[left] at (-7.3, -4) {Pin 5: INJ};
\coordinate (P_CH)  at (-7.25, -6.5);\node[left] at (-7.3, -6.5) {Pin 7: CAN H};
\coordinate (P_CL)  at (-7.25, -7.5);\node[left] at (-7.3, -7.5) {Pin 1: CAN L};


% ==================================================================
% 2. ESP32 (CENTRUM) - GŁÓWNY MODUŁ
% ==================================================================
\draw (6, 0) node[draw, very thick, fill=white, minimum width=6cm, minimum height=13cm, rounded corners] (ESP) {};
\node[above, font=\bfseries\large] at (ESP.north) {ESP32-WROOM-32D};

% Opisy Pinów ESP32
% Lewa strona (Wejścia)
\node[right] at (ESP.west |- 0, 6) {VIN (5V)};
\node[right] at (ESP.west |- 0, 5) {GND};
\node[right] at (ESP.west |- 0, 3) {3V3 (Out)};
\node[right] at (ESP.west |- 0, 1) {GPIO 34 (ADC)};
\node[right] at (ESP.west |- 0, 0) {GPIO 13 (ILL)};
\node[right] at (ESP.west |- 0, -2) {GPIO 26 (SPD)};
\node[right] at (ESP.west |- 0, -4) {GPIO 27 (INJ)};

% Prawa strona (Wyjścia)
\node[left] at (ESP.east |- 0, 6) {GPIO 18 (SCK)};
\node[left] at (ESP.east |- 0, 5) {GPIO 23 (MOSI)};
\node[left] at (ESP.east |- 0, 4) {RES/DC/CS (4,2,5)};
\node[left] at (ESP.east |- 0, -2) {GPIO 16 (RX2)};
\node[left] at (ESP.east |- 0, -3) {GPIO 17 (TX2)};


% ==================================================================
% 3. ZASILANIE (BUCK CONVERTER 12V -> 5V)
% ==================================================================
\node[draw, fill=white, minimum width=3cm, minimum height=2cm, text width=3cm, align=center] (BUCK) at (-2, 6.5) {\textbf{PRZETWORNICA}\\12V-5V};
\node[left, font=\tiny] at (-3.5, 7) {IN+};
\node[left, font=\tiny] at (-3.5, 6) {IN-};
\node[right, font=\tiny] at (-0.5, 7) {OUT+};
\node[right, font=\tiny] at (-0.5, 6) {OUT-};

% Połączenia INPUT (IGN i GND)
\draw (P_IGN) -- (-4, 7) -- (-3.5, 7);
\draw (P_GND) -- (-4, 6) -- (-3.5, 6);

% C1 (Kondensator Wejściowy - Elektrolit)
\draw (-4, 7) to[curved capacitor, l=100$\mu$F (25V), *-*] (-4, 5) node[ground]{};
\node[red, font=\bfseries] at (-3.6, 6.8) {+};

% Połączenia OUTPUT (5V do ESP)
\draw (-0.5, 7) -- (2, 7) -- (2, 6) -- (3, 6); % 5V linia
\draw (-0.5, 6) -- (-0.5, 6) node[ground]{}; % GND Bucka

% C2 (Kondensator Wyjściowy - Elektrolit)
\draw (2, 7) to[curved capacitor, l=100$\mu$F (10V), *-] (2, 5.5) node[ground]{};
\node[red, font=\bfseries] at (2.4, 6.8) {+};

% Masa ESP
\draw (3, 5) -- (2, 5) node[circ]{};


% ==================================================================
% 4. WOLTOMIERZ (ANALOG)
% ==================================================================
% Dzielnik napięcia (R1 = 100k, R2 = 20k)
\draw (P_BAT) -- (-5, 3) to[R, l=100k, -*] (-1, 3) coordinate(DIV);
\draw (DIV) -- (-1, 2) to[R, l=20k] (-1, 1.5) node[ground]{};

% Połączenie do ESP
\draw (DIV) -- (-1, 1) -- (3, 1);

% C3 (Filtr Woltomierza - Elektrolit)
\draw (DIV) -- (0.5, 3) to[curved capacitor, l=2.2$\mu$F, *-] (0.5, 1.5) node[ground]{};
\node[red, font=\bfseries] at (0.9, 2.8) {+};


% ==================================================================
% 5. WEJŚCIA CYFROWE (ZABEZPIECZENIE ZENERA)
% ==================================================================
% Używamy pętli lub makra dla czytelności, ale tu rozpisane ręcznie dla jasności
% Odstępy są duże, żeby teksty się nie zlewały

% ILL (Światła)
\draw (P_ILL) -- (-5, 0) to[R, l=20k, -*] (1, 0) coordinate(Z1);
\draw (Z1) -- (3, 0);
\draw (Z1) to[zzD, l=3V3 Zener, invert, *-] (1, -1.2) node[ground]{};

% SPD (Prędkość)
\draw (P_SPD) -- (-5, -2) to[R, l=20k, -*] (1, -2) coordinate(Z2);
\draw (Z2) -- (3, -2);
\draw (Z2) to[zzD, l=3V3 Zener, invert, *-] (1, -3.2) node[ground]{};

% INJ (Wtrysk)
\draw (P_INJ) -- (-5, -4) to[R, l=20k, -*] (1, -4) coordinate(Z3);
\draw (Z3) -- (3, -4);
\draw (Z3) to[zzD, l=3V3 Zener, invert, *-] (1, -5.2) node[ground]{};


% ==================================================================
% 6. MODUŁY PERYFERYJNE (LCD + CAN) - PRAWA STRONA
% ==================================================================

% LCD Okrągły
\draw (14, 5) node[draw, circle, very thick, minimum size=3.5cm, align=center] (LCD) {\textbf{LCD}\\GC9A01\\(SPI)};

% Moduł CAN
\node[draw, fill=white, minimum width=3cm, minimum height=2.5cm, text width=3cm, align=center] (CAN) at (14, -2.5) {\textbf{MODUŁ CAN}\\SN65HVD230};

% Zasilanie 3.3V (ze stabilizatora ESP)
\draw (3, 3) -- (4, 3) -- (4, 9) -- (10, 9) coordinate(BUS3V3);
\node[above] at (7, 9) {\textbf{Szyna Zasilania 3.3V}};

\draw (BUS3V3) -| (LCD.north); % Do LCD
\draw (BUS3V3) -| (12, -0.5) -- (CAN.north); % Do CAN (omija LCD)

% GND Peryferiów
\draw (LCD.east) -- (16.5, 5) -- (16.5, 4) node[ground]{};
\draw (CAN.east) -- (16.5, -2.5) -- (16.5, -3.5) node[ground]{};

% Magistrala SPI
\draw (9, 6) -- (10, 6) -- (12.25, 6); % SCK
\draw (9, 5) -- (10, 5) -- (12.25, 5); % MOSI
\node[above, font=\small] at (11.5, 6) {SPI Bus};

% UART do CAN
\draw (9, -2) -- (11, -2) |- (12.5, -2.5); % RX
\draw (9, -3) -- (11, -3) |- (12.5, -3.5); % TX

% Wyjścia CAN do wtyczki (Pętla dołem)
\draw (CAN.south) -- (14, -5) coordinate(CAN_OUT);
\draw (CAN_OUT) -- (14, -9) -- (-8, -9) -- (-8, -6.5) -- (P_CH); % CAN H
\draw (-8, -9) -- (-8, -7.5) -- (P_CL); % CAN L
\node[above] at (3, -9) {\textbf{Magistrala CAN (Skrętka)}};


\end{circuitikz}
\end{document}